\section{Results}
\label{sec:results}

\subsection{The system and the upstream ladder}
\label{sec:res:main}

%% 目的：給部署系統的分數與區間，並把測試側的 like-for-like 講清楚（0.5983 對 shipped 自己的 0.6087）。
%%       三個 RQ 的列舉已刪——四個子節標題逐字說著同樣的話。排行榜取證（3000 次分割、0.6225）也刪，
%%       那些數字這一段自己說「不拿來接任何東西」，而 §7.1 才真正量它。
Table~\ref{tab:main} gives the deployed system on val-1000 with bootstrap
intervals, and three ensembles spliced as the binary source into one
frozen downstream pipeline. One interval property matters later: company-clustered resampling
moves the total half-width only from $0.0222$ to $0.0230$, and narrows it on two of the four
fields, so we report both; Section~\ref{sec:discussion} says why that points at the training side
rather than the evaluation side.

The system's one official run read $0.6760$ public and $0.6405$ private on the live
leaderboards; rescored post hoc against the released labels under the schema label set of
Eq.~\eqref{eq:score}, the same predictions give $0.6275$, below both readings, a gap consistent with a different label set
(Section~\ref{sec:discussion}); all three are leaderboard-informed (Section~\ref{sec:setup}). Nor is $0.6275$ the
like-for-like figure: it includes five items
hand-forced to \textsc{Misleading}, of which the labels confirm one, and the validation side
carries no such intervention because the forced identifiers lie outside its range. Withdraw the
confirmed one and test reads $0.5983$ against the shipped build's own $0.6087$ on validation
(Section~\ref{sec:system}), a gap of $0.0104$; Section~\ref{sec:res:price} prices what the five
forced items did.

\begin{table}[t]
  \centering
  \caption{\textbf{Top}: the deployed roster on val-1000, reconstructed as in
  Section~\ref{sec:system}; 95\% percentile bootstrap ($B{=}10{,}000$), ``cluster'' resamples the
  50 companies; intervals condition on the frozen predictions and exclude model-selection
  uncertainty. \textbf{Bottom}: the same frozen downstream with only the binary source changed,
  so the deployed roster's own row reproduces the total above; the 14-model row takes the
  downstream multi-task checkpoints as the binary source.}
  \label{tab:main}
  \footnotesize
  \begin{tabular}{llccc}
    \toprule
    Field & $w_f$ & $M_f$ & 95\% CI (iid) & 95\% CI (cluster) \\
    \midrule
    promise   & 0.20 & 0.8148 & [.7807,.8468] & [.7844,.8459] \\
    evidence  & 0.30 & 0.7083 & [.6706,.7433] & [.6635,.7496] \\
    quality   & 0.35 & 0.3855 & [.3551,.4141] & [.3559,.4137] \\
    timeline  & 0.15 & 0.6607 & [.6089,.7032] & [.6004,.7053] \\
    total $S$ &      & \textbf{0.6095} & [.5864,.6308] & [.5854,.6313] \\
    \midrule
    \multicolumn{5}{l}{\textit{binary source (downstream frozen)}} \\
    \multicolumn{2}{l}{\quad diluted 7-model vote}  & \multicolumn{3}{l}{$0.6053$} \\
    \multicolumn{2}{l}{\quad \textsc{strong4} (deployed)} & \multicolumn{3}{l}{$0.6095$} \\
    \multicolumn{2}{l}{\quad 14-model ensemble}      & \multicolumn{3}{l}{$0.6108$} \\
    \bottomrule
  \end{tabular}
\end{table}

\subsection{RQ1: what is the promise gate worth?}
\label{sec:res:weight}

%% 目的：RQ1：在部署操作點算出 promise 的有效權重 0.5216（另一方向 0.4268），並按目的地欄位拆解，
%%       得出 62% 的貢獻收在 promise 之外。CLaC 那一段移到 §2（那裡已經帶著 0.2330/0.2575 講完），
%%       段末重述命題的那句也刪掉——同一句話這一節已經是第四次。
At the deployed operating point ($\rho = 0.9828$, $\varphi_T = 0.4385$) the counted model answers
RQ1 directly: marginal gate recall moves the total at
$\mathrm{d}S/\mathrm{d}\rho = +0.6078$ and the promise field at
$\mathrm{d}M_P/\mathrm{d}\rho = +1.1652$, so promise's effective weight is
$W_{\mathrm{eff}}(P) = 0.6078/1.1652 = \mathbf{0.5216}$ (bootstrap $95\%$ CI
$0.485$--$0.542$; company halves $0.513$/$0.524$) against a nominal $0.20$, a factor of $2.6$. The specificity
direction gives $0.4268$ ($0.403$--$0.450$), $2.1\times$ nominal: the published $0.20$ locates
neither. By destination field, $\mathrm{d}S/\mathrm{d}\rho$ splits into promise $+0.2330$,
evidence $+0.2575$, quality $+0.0720$, timeline $+0.0452$: \textbf{$61.7\%$ of the marginal
contribution of a better promise decision is collected outside the promise field}, the largest
channel being the evidence field's \emph{scored} N/A class (Section~\ref{sec:metric}).
Recounted on the released test labels at the shipped submission's own operating point, all three
reappear: $0.5267$, $0.4374$, $62.0\%$ outside.

%% 目的：讓 decomposition 可否證——一次有限的改善應該落在導數說它會落的地方。實測只對一半：
%%       導數說 61.7% 收在 promise 之外，完整 gold-splice 說 51.9%。方向與通道排序成立，份額不成立。
The decomposition is falsifiable against a finite move, and it is only partly right:
gold-splicing promise the whole way (Section~\ref{sec:res:oracle}) puts $51.9\%$ of the gain
outside the promise field against the derivative's $61.7\%$. The direction survives; the leading
channel and the share do not, and the gap is the curvature between a marginal and a total
change.

\subsection{RQ2: is the cascade model predictive?}
\label{sec:res:swap}

%% 目的：RQ2：換上游、下游全凍，看模型能不能預測沒被動過的兩個欄位。誠實揭露一半是恆等式，
%%       並給出「有條件的否定」——平均值 0.0020、v1 0.0014、v2 0.0023，越低越好，我們部署的那個沒贏。
Because the model has no tuned parameters (Section~\ref{sec:identify}), the swap asks what its one
empirical assumption costs: that a downstream head's accuracy given an open gate does not depend
on which gate opened it. Holding that fixed, we swap only the binary source across the 42
independently trained single models of Section~\ref{sec:setup}, which form ten backbone families;
\textsc{strong4} is an ensemble over four of them and is not itself among the 42. Gate parameters
are recounted per source; nothing downstream is retrained or re-thresholded, so all timeline and
quality movement is cascade by construction. Observed promise spans $.769$--$.815$, and the two
unmodified fields move with it: timeline $.651$--$.661$, quality $.368$--$.396$.

Half the composite is not at risk here and we say so: given a source's gate counts, the promise
and evidence macros are the observed tables rewritten, reproduced to $10^{-16}$ by arithmetic
rather than by prediction. What the swap predicts is $U = 0.35\,M_Q + 0.15\,M_T$, spanning
$0.0092$ across the 42 sources. Against a leave-one-out mean of $U$ --- the fair comparator once
those identities are granted to it too --- mean absolute error is $0.0020$, the class-independent
variant reaches $0.0014$ and the class-conditional one $0.0023$, lower being better. RQ2 therefore
gets a qualified no for the variant we deploy, and the margin by which the other one wins is not
resolvable from ten families on one split. We keep $\gamma_c$ for identification, not accuracy.

%% 目的：參數本身才是發現：級聯傷害不均勻，而且在 test 上重數也站得住（水準差在 0.05 以內）。
The parameter, not the prediction, is the finding, and it is counted rather than fitted: cascade
damage is not uniform. The gate closes on $16.8\%$ of gold \textit{Not Clear} items but only
$4.1\%$ of \textit{Clear} ones ($\gamma_Q = 0.832$ and $0.959$); $\gamma_T$ runs from $0.974$ on
\textit{already} to $1.000$ on $\le$2y. In quality the gate closes hardest on the class the
downstream head is already weakest on (Section~\ref{sec:errors}), which macro-averaging weights as
heavily as the strong one; timeline places its closures the other way, on \texttt{already}, and
loses less throughout. Because one system is scored on both splits this can be checked: on the
released test labels $\gamma_Q$ is $0.788$ on \textit{Not Clear} against $0.943$, and $\gamma_T$ is $0.948$ on
\textit{already} --- ordering and levels both hold, within $0.05$ on every one.

\subsection{RQ3a: how much headroom, and where?}
\label{sec:res:oracle}

%% 目的：用 gold-splice 的 2x2 oracle 把上限量化（+0.1565 = direct 0.1246 + cascade 0.0319），
%%       原本花了 74 字宣告一個反對意見再回答它、再用文字重述一次答案，範圍本身就是答案。
%%       「被救回的資料是不是比較難」那四句（timeline 那半邊結論是「無法判定」）壓成一個子句。
To bound the cascade we splice gold labels into the gate in a $2\times2$ design, downstream
unchanged (Table~\ref{tab:oracle}). A perfect binary contributes $+0.1565$: $+0.1246$
\emph{direct}, in promise and evidence themselves, and $+0.0319$ ($20.4\%$) \emph{cascade}, in
the timeline and quality fields that no module modified. Promise alone accounts for $0.0770$ of
that total and a third of the cascade term, so the cascade term belongs to the binary source as a
whole rather than to promise. Under a perfect gate the downstream modules cap
at $0.6905$ on timeline and $0.4640$ on quality --- ceilings no upstream effort moves
(Section~\ref{sec:errors}). Rescued rows are harder in quality ($0.675$ right, $n{=}40$, against $0.839$ of
the already-open ones), so the $+0.0319$ is a floor there.

\begin{table}[t]
\centering\footnotesize
\caption{Top: gold-splice oracle on validation, direct versus cascade points. Bottom: what one
repaired error contributes, in $10^{-4}$ points, with how many are available. Each price is the mean gain
from repairing all $n$ such rows to gold, promise repairs carrying the evidence they gate; sums
across rows are first-order only.}
\label{tab:oracle}
\begin{tabular}{@{}lccccc@{}}
\toprule
Gate / scored source & $S$ & P & E & T & Q \\
\midrule
B\; model gate, model scored & 0.6095 &.8148 &.7083 &.6607 &.3855 \\
O$_P$\; gold promise & 0.6864 & 1.000 &.8067 &.6905 &.4024 \\
O$_E$\; gold evidence & 0.6761 &.8148 &.8708 &.6607 &.4364 \\
O$_{PE}$\; both gold & \textbf{0.7660} & 1.000 & 1.000 &.6905 &.4640 \\
\midrule
\multicolumn{6}{@{}l@{}}{headroom $+0.1565 =$ direct $+0.1246 +$ cascade $+0.0319$ (20.4\%)}\\
\midrule
\multicolumn{4}{@{}l}{Repair} & points & $n$ \\
\midrule
\multicolumn{4}{@{}l}{one true \textsc{Misleading} caught} & \textbf{63.9} & 1 \\
\multicolumn{4}{@{}l}{quality conditional fix} & 7.0 & 101 \\
\multicolumn{4}{@{}l}{promise false negative} & 8.5 & 14 \\
\multicolumn{4}{@{}l}{evidence error, gate open} & 6.6 & 106 \\
\multicolumn{4}{@{}l}{promise false positive} & 8.8 & 82 \\
\multicolumn{4}{@{}l}{timeline conditional fix} & 1.7 & 266 \\
\bottomrule
\end{tabular}
\end{table}

\subsection{RQ3b: what is one repair worth, or one rare class?}
\label{sec:res:price}

%% 目的：RQ3b：把每一種可修錯誤換算成總分單位，凸顯 Misleading 那一列遠超其他、且我們一分未得。
%%       原本「138 次修復」旁邊就是 n=14，讀者會想去比一個不存在的量，補上「而實際只有 14 個」。
The same model expresses every repairable error in total-score units (Table~\ref{tab:oracle}),
computable before any submission from the field weights, label distribution and operating point
alone. By Lemma~1 the \textit{Misleading} class is worth $0.1167$ whatever its support --- as much
as $138$ promise false-negative repairs would return, against the 14 such errors that exist ---
which exceeds the whole cascade headroom ($0.0319$) and rivals the direct-binary headroom
($0.1246$). On validation we score $F_1{=}0$ there, so none of it is recovered. What a single
\emph{catch} returns depends on what the system already emits: the repick places 35 validation
items in that class, none correct, so adding the one true item ($k{=}36$) returns $0.0064$ ---
$0.0063$ of it \textsc{Misleading} at $F_{1,c}=2/37$. The slot is large because the
class is unclaimed; the catch is small because the slot is already spent badly.

\leadin{That contribution is observable on our own submission.} The shipped file forces five items
to \textsc{Misleading} and the released labels confirm exactly one. Withdrawing the four wrong
ones, one at a time, traces the curve on the artefact we actually submitted:

{\centering\footnotesize
\begin{tabular}{@{}lcccccc@{}}
\toprule
items forced \textsc{Misleading} & 4$^{\dagger}$ & 5 & 4 & 3 & 2 & 1 \\
\midrule
test quality macro & .3676 & .4511 & .4629 & .4786 & .5010 & \textbf{.5342} \\
total $S$ & .5983 & \textbf{.6275} & .6317 & .6372 & .6450 & .6566 \\
\bottomrule
\end{tabular}\par}
\noindent{\footnotesize $^{\dagger}$the four false positives alone, with the confirmed item
withdrawn. Column two is the submitted file.}

Withdrawing false positives raises the field monotonically, $0.0831$ in all; the set that keeps
only the four wrong guesses scores $.3676$, which is what forcing nothing at all scores ($.3675$)
--- four confident placements on the rare class are worth effectively nothing. Every column is
Lemma~1 arithmetic: with one of the three gold instances found, precision $1/k$ over $k$ forced
items gives the class $F_1 = 2/(3+k)$. The two directions are not interchangeable: at the
submitted $k{=}5$ a withdrawal adds $0.0119$ to the field and a further catch adds $0.0833$, a
factor of seven. Against the $.3675$ of forcing nothing, the one correct placement is worth
$+0.0292$ of the total --- more than the $+0.0042$ our ensemble search returned
(Table~\ref{tab:main}). That
is a statement about the scorer, not about detection (Section~\ref{sec:errors:mis}).

%% 目的：裁決。§1 承諾「解決這個衝突才是本文要做的事」，但原本 Results 把兩邊都定完價就停了，
%%       讀者手上握著兩個數字沒有結論——而最大的那個是 0.1167，他會自己推論「去追 Misleading」。
%%       這一段是全篇唯一讓 0.1167、0.0319、0.5216 並排出現並給出判決的地方。
\leadin{The conflict, resolved by these numbers.} Lemma~1 makes the rare class tied for the largest
scored \emph{class} on the board, $0.1167$ --- above either field's own oracle gain (promise
$0.0770$, evidence $0.0666$), though below the two-field binary headroom of $0.1565$, which is
neither one class nor one repair. Eq.~\eqref{eq:weff} makes the gate worth $0.5216$ per point of
its own macro, $2.6$ times its published weight. The
two prices, in different units, do not settle it, because a price is not an invoice: of the rare
class's $0.1167$ we
collected nothing on validation and $0.0292$ on test, from one placement out of five guesses,
while the gate's gains were reachable by ordinary subset selection. For our system, only upstream
effort showed any gain on validation; at supports of 1 and 3 we draw no general rule for
the rare class (Section~\ref{sec:errors:mis}).
